\begin{afterword}
\summaryhead

The author recalls a panel on whether science and religion are compatible. A pagan panelist recited at length a creation myth in which a primordial cow licks a god into existence and the god's descendants build the Earth from a giant's corpse. The author is sure that she knew it was scientifically absurd. She held that religion and science were compatible, and even explained the myth by the Vikings' landscape, yet said with odd pride that this was what she ``believed.''

The author asks what kind of belief this was. It was not Dennett's ``belief in belief,'' since she showed no fanatical need to reassure herself. It was not ``religious profession'' either, since anyone professing a belief to satisfy others or themselves would pretend more convincingly. The author concludes that she was cheering for paganism, like holding a GO BLUES banner at a football game. Her pride was like marching naked in a pride parade: the outrageousness was the point, and a more plausible version would have been like putting on clothes.

\responsehead

The post draws a real distinction. Saying a sentence to assert something and saying it to show which side you are on are different acts, and the banner at a football game makes the difference easy to see. ``Belief as Attire,'' posted the same day, builds on it, with a metaphor Robin Hanson offered in a comment on this post.

The case the post uses to introduce the distinction is weaker than the distinction. There is one panelist, recalled by the author, whose state of mind is read from her manner. The author is candid that the manner was hard to name: ``pride? Self-satisfaction? A deliberate flaunting of herself?'' From that impression the post draws three claims about her mind: that she knew the myth was scientifically outrageous, that she was not trying to be convincing, and that she believed she ``couldn't be arrested or criticized.'' The first is introduced with ``clearly,'' twice. Each is based on the impression.

The two things she is reported to have said suggest a different reading. She held that religion and science were compatible, and she explained the myth by the land the Vikings lived in. Both fit a reading the post does not consider: that she took the myth as symbolic, as most modern followers of Norse paganism take their cosmology, according to the academic studies Wikipedia summarizes (Wikipedia adds that most of them believe the gods themselves exist). On that reading she had nothing to pretend and nothing to flaunt. She was telling a story she valued without taking it as history. The post sets ``believed'' in quotation marks, as if the word were being used oddly, but does not report asking her what she meant by it.

As with ``Belief as Attire'' (see the notes there), the post gives no way to tell a cheer from a belief, so the label depends on the observer's impression. The one general claim it adds, that ``most people are honest enough'' to repeat a spoken religious statement to themselves, comes without evidence.

In short: a useful distinction between asserting and cheering, drawn from one recalled panelist whose mind the post reads from her manner, without considering the symbolic reading that her own reported words suggest.
\end{afterword}
